\section{Deployment, Lessons Learned, and Limitations}
\label{sec:lessons}

\paragraph{Deployment.}
\process{} runs as a command-line pipeline with resumable state: an analyst seeds a mineral, lets discovery run overnight, inspects the \tool{} and its references, edits the material set, and resumes. Reference content is cached and indexed for semantic search, so later iterations and later minerals reuse evidence already fetched. Three design decisions followed from the analysts' stated needs rather than from research convenience: every vertex is typed with an HS-6 code so the \tool{} can be joined to trade data (Section~\ref{sec:mekh}) -- though, as Section~\ref{sec:experiments} shows, that typing is not guaranteed correct; every hyperedge carries scored references so an analyst can audit it (Section~\ref{sec:gates}); and the purge step removes rather than flags unsupported processes, because the analysts preferred a smaller trusted graph to a larger annotated one. Beyond the four \tool{}s evaluated in Section~\ref{sec:experiments}, the same pipeline and generator model (\GenModel{}) were also run once for the sponsor as a demonstration on helium, a gas whose supply chain runs through isotope separation (helium-3, lithium-6, tritium); that run was a demonstration, not an evaluated \tool{}, and we report no numbers from it here.

\paragraph{What goes wrong.}
Three failure modes dominate. (1)~\emph{Code drift}: an agent identifies the right material but assigns a neighbouring HS-6 subheading, or cannot code it at all and records it as an untyped plain string (Section~\ref{sec:gates}) or under a code no vertex represents. This is the dominant RQ1 failure -- HS codes are judged correct for only \JudgeHSGalliumClaude{}--\JudgeHSCobaltClaude{} of hyperedges even though the underlying chemistry is judged plausible almost every time (Section~\ref{sec:experiments}) -- and it produces concrete recall misses: gallium arsenide coded as a gallium-metal article, trimethylgallium given a non-existent HS code, and germanium tetrachloride coded inconsistently across processes. Constraining code assignment to the subheadings of a heading already present in the graph, and checking assigned codes against the official nomenclature rather than trusting the search tool's top result, would target both symptoms. (2)~\emph{Narrative-only inputs and outputs}: a process description mentions a reagent or product in prose that is never added to the structured precursor or product list, so it is invisible to any analysis keyed on those lists. The clearest case is electrowon cobalt metal, which appears in a process narrative but not in that process's product list. A schema field that forces every reagent or product named in the narrative to be resolved to a material, typed or not, would recover such cases. (3)~\emph{Coverage gaps in industrially central routes}: some routes never appear at all, not from a coding failure but because no agent proposed them in any iteration -- MOCVD growth of gallium nitride from trimethylgallium, the lithium-cobalt-oxide cathode chemistry, and zone refining of elemental germanium are all missing from the corresponding \tool{}. Seeding discovery with a checklist of known end uses per mineral, rather than relying only on emergent proposals from the classification and manufacturing agents, would be one way to close such gaps.

\paragraph{Limitations.}
The LLM judges are a proxy for expert review, not a substitute for it, and we report their verdicts as given: \JudgeAName{} and \JudgeBName{} agree only partially (Table~\ref{tab:judge}), and we do not adjudicate cases where a judge's rationale and its verdict disagree. Evidence grounding is measured by word-boundary, alias-aware string matching, which under-counts paraphrase -- a reference can describe the right process without naming a material as the \tool{} names it. Canonical-route recall depends on small, author-curated route sets (\RoutesNCobalt{}--\RoutesNBoron{} routes per mineral, drawn from public sources), so a single route moves reported recall by double digits. HS-6 aggregation merges materials of different purity into one vertex, and iteration is bounded at five rounds, so routes reachable only through longer chains are never discovered. All four \tool{}s evaluated here are in one domain; that the architecture would generalize beyond critical minerals is a design claim, not something we measure. The criticality and disruption-cascade analyses of Section~\ref{sec:application} are prior results computed on top of the \tool{}, not additional evaluation of it, and inherit every limitation above.

\paragraph{Ethical and responsible-use considerations.}
The pipeline reads only public sources and produces analyses of public trade data; it does not model firms or individuals. Its outputs inform policy analysis and are audited by analysts before use, which the reference requirement is designed to support. AAMAS 2027 policy requires disclosure of generative AI used in a paper's methodology; here that use has three parts, disclosed with the specific models involved in Section~\ref{sec:experiments}: the agents that are this paper's object of study (\GenModel{}), the LLM judges used as an evaluation proxy (\JudgeAName{} and \JudgeBName{}), and editorial polishing of the text. All analyses and claims are the authors' responsibility.

\paragraph{Availability.}
Code is available in an anonymized repository (\url{ANONYMIZED-REPOSITORY-LINK}); agent prompts, evaluation scripts, judge outputs, the curated route and USGS-code files, and the \tool{} state files are in the supplementary material and will be archived on Zenodo on publication.
